% Results: the ten-gate decision surface (centerpiece), the executed control
% behaviors, the 32-iteration sweep and four-entrant tournament, the featured
% 1000-row sensor stream, the 60-second timeline and safety surface, and
% determinism. Dense numbers live in body-width tables; figures 03 and 11 are
% placed here and referenced with \autoref.

The results lead with the decision surface, because it is the operational proof
of the thesis, and then move outward to the behaviors that feed it, the
integration layer that repeats it, the featured artifacts that stress it, and the
determinism that lets every number stand as reproducible evidence. Throughout,
the dense values are carried in the tables and the figures and are referenced
rather than restated.

\subsection{The ten-gate decision surface}

The assurance layer was run over five cases that together exercise the full
accept, block, and escalate surface. In the nominal case the observed metric for
each gate equalled its independent reference, every verification, validation, and
uncertainty dimension passed, a human review was recorded, and all ten gates
returned accept; this is the only path on which the composite score is even
reported. Four adversarial cases then proved that the failing paths actually
fail. A catastrophe hard-predicate failure blocked the vascular no-fly gate even
though every graded dimension passed, because the count of hard-stop volume
breaches was non-zero. A verification fraction of 0.80 blocked the perception
gate, since four of five structural checks passing is still not the exact 1.0 the
gate demands. A validation collapse blocked the balance gate, where driving the
observed zero-moment-point margin 50 percent off its reference dropped agreement
to 0.00 and breached the 0.03 relative-error bound. And a dispersion of 0.163,
above the 0.10 bound, escalated the finger-force gate while verification and
validation both passed, which is the quantification dimension doing its job and
defaulting to hand-back-to-human. \autoref{tab:decisionsurface} summarises the
five cases and the dimension each exercised.

\begin{table}[H]
\caption{The five executed VVUQ decision cases (from
\texttt{papers/VVUQ-02/execution/03-vvuq/}); three distinct block mechanisms, one
escalate, and one full-accept sweep.}
\label{tab:decisionsurface}
\centering
\footnotesize
\begin{tabularx}{\textwidth}{L{1cm} L{3.4cm} Y L{1.9cm}}
\toprule
\textbf{Case} & \textbf{Gate} & \textbf{Dimension exercised} & \textbf{Decision} \\
\midrule
A & all ten gates & nominal, observed equals reference & ten ACCEPT \\
B & 06 vascular-no-fly & hard catastrophe predicate (hard-stop count) & BLOCK \\
C & 08 perception-scene & verification fraction 0.80 below 1.0 & BLOCK \\
D & 03 whole-body-balance & validation agreement 0.00 and rel err 0.500 & BLOCK \\
E & 02 dexterous-finger-force & uncertainty CV 0.163 above 0.10 & ESCALATE \\
\bottomrule
\end{tabularx}
\end{table}

The reason the surface behaves this way is that the ten gate thresholds tighten
with model risk toward the three immediate-catastrophe gates, as
\autoref{fig:forest} makes visible and \autoref{tab:thresholds} states exactly.
Every gate requires a verification fraction of exactly 1.0; what varies is the
validation agreement, the maximum relative error, and the dispersion bound, all
of which are strictest for the vascular no-fly, shared-operating-room collision,
and fault e-stop gates. The harness blocks or escalates on the first failing
dimension and accepts only when verification, validation, quantification, and any
hard predicate all hold, so the decision is categorical, not a score to be traded
off.

\begin{table}[H]
\caption{The ten VVUQ gate thresholds (from
\texttt{papers/VVUQ-02/codegen/config/vvuq\_thresholds.yaml}). Every gate also
requires a verification fraction equal to 1.0; the three catastrophe gates carry
an extra hard predicate.}
\label{tab:thresholds}
\centering
\footnotesize
\begin{tabularx}{\textwidth}{L{0.5cm} Y L{1.9cm} L{1.7cm} L{1.5cm} L{1.5cm}}
\toprule
\textbf{\#} & \textbf{Gate} & \textbf{Agreement ($\geq$)} & \textbf{Max rel err ($\leq$)} & \textbf{CV bound ($\leq$)} & \textbf{Catastrophe} \\
\midrule
01 & bimanual-handeye-servo & 0.97 & 0.050 & 0.08 & no \\
02 & dexterous-finger-force & 0.95 & 0.050 & 0.10 & no \\
03 & whole-body-balance & 0.98 & 0.030 & 0.06 & no \\
04 & autonomous-plan-correctness & 0.95 & 0.050 & 0.10 & no \\
05 & instrument-grasp-handover & 0.96 & 0.050 & 0.10 & no \\
06 & vascular-no-fly-hand & 1.00 & 0.010 & 0.05 & yes \\
07 & bimanual-suturing-anastomosis & 0.96 & 0.050 & 0.08 & no \\
08 & perception-scene-understanding & 0.95 & 0.050 & 0.10 & no \\
09 & shared-or-human-collision & 1.00 & 0.020 & 0.06 & yes \\
10 & fault-estop-graceful-degradation & 1.00 & 0.020 & 0.05 & yes \\
\bottomrule
\end{tabularx}
\end{table}

\subsection{The executed control behaviors}

The candidates the gates judge were produced by the six behavior groups, and each
behaved as its gate expects. \autoref{tab:behaviors} collects the headline result
from each group: the autonomy intent and compile reached a phase-step concordance
of 1.000 against the annotated reference plan, the right seven-axis arm solved to
a tip pose in the world frame, scene segmentation degraded smoothly under
occlusion with all five structures still usable after near-infrared, ultrasound,
and bile-channel fusion, the fingertip force tracked under its soft cap with the
summed two-hand force inside the hard cap, the balance margin held well inside the
support polygon across the disturbance range, and all three anastomoses held ring
tension in band for every stitch. This subsection is intentionally tight, because
these behaviors are the cheap, fast half of the workflow; the substantial work is
the gating already shown, and the integration that follows.

\begin{table}[H]
\caption{Executed control-behavior headlines (from
\texttt{papers/VVUQ-02/execution/02-pipeline/}); each group is the candidate its
gate then judges.}
\label{tab:behaviors}
\centering
\footnotesize
\begin{tabularx}{\textwidth}{L{4.2cm} Y}
\toprule
\textbf{Behavior group (gate)} & \textbf{Headline result} \\
\midrule
Autonomy intent and compile (04) & phase-step concordance 1.000 for phases 1 and 5 \\
Forward kinematics (01) & right arm tip pose solved from the seven-axis chain \\
Perception and fusion (08) & mean Dice 1.000 to 0.819 over occlusion 0.00 to 0.30; five of five structures usable \\
Hands, force and grasp (02, 05) & fingertip 1.84 N under soft cap; summed 5.0 N within hard cap; grasp and handover succeed \\
Whole-body balance (03) & margin 130.0 to 62.3 mm over a 0 to 40 N push; stable, no e-stop \\
Suturing and anastomosis (07) & three anastomoses, RMSE 0.0033 to 0.0047 N, in-band 1.0 \\
\bottomrule
\end{tabularx}
\end{table}

\subsection{The 32-iteration sweep and the four-entrant tournament}

The integration layer repeats the whole pipeline under a deterministic
32-iteration Latin hypercube and then stages a four-entrant comparison. All 32 of
32 iterations cleared every one of the ten gates, and the gated composite spanned
a minimum of 93.417, a maximum of 93.715, and a mean of 93.562; the composite is
reported only because every gate accepted, since the gating overlay withholds the
headline number on any block or escalate, which is the thesis made mechanical.
The comparison is the first of the two featured generated artifacts. It is a
1790-line structured results file of 32 iterations times four rounds, so 128
round verdicts, with 100 percent of the robot-involving verdicts carrying the
simulation-against-simulation caveat and every human-versus-robot round
additionally carrying the structural time-dimension caveat.
\autoref{tab:leaderboard} reads the leaderboard directly from that file.

\begin{table}[H]
\caption{Four-entrant tournament leaderboard (from
\texttt{papers/VVUQ-02/codegen/results/comparison.json}); frozen-weight composite,
simulation against simulation.}
\label{tab:leaderboard}
\centering
\footnotesize
\begin{tabularx}{\textwidth}{L{0.9cm} Y L{2.3cm} L{1.5cm} L{1.7cm}}
\toprule
\textbf{Rank} & \textbf{Entrant} & \textbf{Composite mean} & \textbf{Win rate} & \textbf{Total wins} \\
\midrule
1 & PancreSpeed 1.0 (8-arm cart) & 93.782 & 0.875 & 56 \\
2 & H2-Surgical 1.0 (humanoid) & 93.334 & 0.75 & 72 \\
3 & da Vinci-successor 2030 (teleop) & 83.970 & 0.0 & 0 \\
4 & Dutch human baseline (2025) & 67.885 & 0.0 & 0 \\
\bottomrule
\end{tabularx}
\end{table}

The single mobile humanoid placed second, behind the eight-arm PancreSpeed cart
by under half a composite point, and the reason is structural rather than a defect
of the humanoid. The composite is weighted toward throughput-sensitive
components, and eight arms working in parallel complete the same eight-phase
procedure with more simultaneous action than two hands can, so the cart edges
ahead on the time-weighted score and wins the head-to-head round 24 of 32 times,
which lifts its win rate to 0.875. The humanoid nonetheless accumulates more total
wins, 72 against 56, because it appears in three of the four rounds while the cart
appears in two, and the win rate normalizes for those appearances. This ordering
is exactly why the comparison is informative: where the prior PDAC paper featured
the multi-arm cart as the lone fast baseline, the present study introduces the
single mobile humanoid against it and finds the humanoid lands within striking
distance of the parallel-arm platform while concentrating all risk into two hands,
which is the higher-assurance case this paper is built around. The
generation-versus-assurance point is explicit here too: producing a 1790-line
structured tournament is fast and deterministic, while judging it, the caveat
coverage, the win accounting, and the dispersion, is the assurance work.

\subsection{The featured humanoid sensor stream}

The second featured artifact is the positional sensor stream, and it is the
substantial code-generation feat on the data side. The committed stream is 1000
unique rows across 27 columns, 500 command ticks for each of the two hands,
carrying seven arm-joint channels, five fingertip force channels, and three
end-effector position channels per hand, plus the safety channels for vessel
proximity, ring tension, support-polygon margin, balance state, forces, e-stop,
and collision. Every whole row and every per-row arm-angle, finger-force, and
end-effector payload is distinct, with zero duplicated rows and no constant-padded
positional column, and the sample is the first 50 ms of phase 1 at the 10 kHz
command rate. \autoref{fig:bands} plots three of its safety-relevant channels
against the bounds their gates enforce, and \autoref{tab:bands} gives the
per-channel statistics. The point the subsection makes is the thesis in miniature
on the data side: generating a thousand non-repetitive rows of two hands, ten
fingers, and arm channels is cheap, while verifying their uniqueness and their
safety-channel ranges is the assurance work.

\begin{table}[H]
\caption{Per-channel statistics of the featured 1000-row stream (from
\texttt{sample\_h2\_sensor.csv} and execution \S05), with the governing gate
bounds.}
\label{tab:bands}
\centering
\footnotesize
\begin{tabularx}{\textwidth}{Y L{1.0cm} L{1.0cm} L{1.0cm} L{0.95cm} L{2.5cm} L{1.1cm}}
\toprule
\textbf{Channel (unit)} & \textbf{Min} & \textbf{Max} & \textbf{Mean} & \textbf{SD} & \textbf{Bound(s)} & \textbf{Gate} \\
\midrule
Bimanual cumulative force (N) & 4.366 & 6.716 & 5.507 & 0.388 & soft 5.0, hard 6.0 & 02 \\
Support-polygon margin (mm) & 38.005 & 42.992 & 40.458 & 1.440 & floor 8.0 & 03 \\
Vessel proximity (mm) & 6.018 & 9.999 & 7.983 & 1.155 & no-fly 6.0, hard-stop 2.0 & 06 \\
\bottomrule
\end{tabularx}
\end{table}

\subsection{The 60-second timeline and the safety surface}

The deployment reference runs the eight-phase 60-second Whipple, from the Kocher
exploration in phase 1 to closure in phase 8, and exercises the three
immediate-catastrophe safety surfaces at their boundaries. The vascular no-fly
gate layers nested keep-out radii around each active vessel and reduces velocity
and finally e-stops as the fingertip approaches, producing zero hard-stop
violations on a safe path, which is the very hard predicate that blocked case B.
The shared-operating-room collision gate drives a four-state machine from full
speed when clear, through half speed in proximity and a crawl on contact, to an
e-stop when the human clearance reaches the floor. The fault gate passes a clean
snapshot, safe-parks a single balance fault within the 3 ms e-stop budget,
hands back on a vision dropout, and escalates to hand-back on any multi-fault.
These bounds are the same ones shown in \autoref{fig:bands}, so the subsection
stays brief and defers to the figure rather than repeating the numbers.

\subsection{Determinism}

The results close where the methods began, on reproducibility. The thousand-row
sensor stream regenerated byte-for-byte, the four-entrant comparison regenerated
as an identical semantic match, and the sweep index regenerated identically, all
from seed 20260525. That byte-level determinism is what lets every number in this
section stand as reproducible evidence a reader can re-derive, rather than a
summary that must be trusted, and it is the property the standards-anchored
assurance argument in the Discussion depends on.

\begin{figure}[H]
   \centering
   \IfFileExists{Images/03-ten-gate-threshold-forest.png}{%
     \hspace{-0.95cm}%
     \includegraphics[width=1.1\linewidth, trim={0 0 0 0}, clip]{Images/03-ten-gate-threshold-forest.png}%
     % \vspace{-0.25cm}
   }{%
     \setlength{\fboxsep}{10pt}%
     \fbox{\parbox[c][6cm][c]{0.9\linewidth}{\centering\itshape Figure placeholder.
     The author adds \nolinkurl{Images/03-ten-gate-threshold-forest.png}, rendered by
     \nolinkurl{papers/VVUQ-02/imagegen/03-ten-gate-threshold-forest/03-ten-gate-threshold-forest.py}.}}%
   }
    \caption{Forest plot of the ten VVUQ gate thresholds, with bounds tightening toward the three immediate-catastrophe gates.}
    \label{fig:forest}
\end{figure}

\begin{figure}[H]
   \centering
   \IfFileExists{Images/11-sensor-stream-safety-bands.png}{%
     \hspace{-0.95cm}%
     \includegraphics[width=1.1\linewidth, trim={0 0 0 0}, clip]{Images/11-sensor-stream-safety-bands.png}%
     % \vspace{-0.25cm}
   }{%
     \setlength{\fboxsep}{10pt}%
     \fbox{\parbox[c][6cm][c]{0.9\linewidth}{\centering\itshape Figure placeholder.
     The author adds \nolinkurl{Images/11-sensor-stream-safety-bands.png}, rendered by
     \nolinkurl{papers/VVUQ-02/imagegen/11-sensor-stream-safety-bands/11-sensor-stream-safety-bands.py}.}}%
   }
    \caption{Live safety margins across the featured 1000-row humanoid sensor stream with mean and uncertainty bands against the governing gate bounds.}
    \label{fig:bands}
\end{figure}
